\section{What a refusal would have to do}
\label{sec:17.1}
Three things separate a worthwhile refusal from what a keyword filter produces. A request arrives: resolve this description of a person into a location. The same capability finds a hostage or a missing child, and nothing in the request says which. If something declines, it has to do all three:

\begin{itemize}
\item \emph{It reaches new cases.} It has to reach the next request, which won't use the same words. A vehicle, a habitual route, and a time of day aren't a description of a person, and together they resolve to the same place. They needn't arrive together, either: the vehicle today, the route tomorrow, the hour next week. So this condition applies to sequences of requests as well as single ones. A refusal tied to the wording fails it.
\item \emph{It lifts when the reason is gone, and only then.} A refusal \emph{lifts} when it stops applying and the system proceeds, and it has to lift when the reason is gone: a parent coordinating a search is asking for the same act. And the fact that lifts it usually isn't in the request; someone outside has to establish it. A refusal that can't lift, or lifts on the requester's say-so, fails it.
\item \emph{It holds under pressure: pressure on the refuser doesn't change the answer.} It weighs what the act would do against what refusing would leave undone, and, within the range of costs the arrangements around it permit, it doesn't give way because refusing has become expensive for whoever refuses.
\end{itemize}

A worthwhile refusal weighs the harm of acting against the harm of refusing. It sees the cost to the refuser, because a refusal that can't see what it costs can't be shown to have held against it, and that cost doesn't move the outcome. Raise the cost to the refuser with nothing else changed, and the refusal stays put. Supply a corroborated fact that raises the harm of refusing, and it moves. No learner, biological or artificial, is known to hold a priority of that kind without limit; the nearest thing in people, the sacred values that resist trade-offs \autocite{ginges2007sacred}, trade under some framings, which is why the floor's absolute terms don't rest on a learner alone: institutions or a mechanism uphold them too, where any remain. So the range over which a refuser is asked to hold is set from outside, by the arrangements around it, and a refusal is tested over that range. What would let it hold past that range is integrity; what makes holding survivable at all is the room it is in. A ranking across kinds of wrong meets the same limit: a learner formed toward it approximates it, and where a deployment needs the ranking to hold without exception, that part is written as a floor term and held outside the system.
